\begin{solapartat}
\punts{0,1} L'equació del moviment serà de la forma: $y(t) = A\cos(\omega t + \phi)$

\punts{0,2} Segons l'enunciat, el desplaçament total del vaixell és de 2~m i, per tant, l'amplitud del moviment és A=1~m i $\dfrac{T}{2} = 6,28\ \mathrm{s} \Rightarrow T = 12,56\ \mathrm{s}$ \ (*)

\punts{0,2} La freqüència angular (pulsació) és:
\[ \omega = \frac{2\pi}{T} = \frac{2\pi}{12,56} = 0,50\ \mathrm{rad\,s^{-1}} \]

\punts{0,2} Segons l'enunciat, y(t=0) = A = 1~m

\punts{0,3} Així, l'equació del moviment del vaixell és: $y(t) = 1\cdot\cos(0,50t)\ \ (\mathrm{m})$ o $y(t) = 1\cdot\sin\left(0,50t + \dfrac{\pi}{2}\right)\ \ (\mathrm{m})$
\end{solapartat}

\begin{solapartat}
\punts{1}

$\mathit{Si}\ y = A$ en t = 0 $\left\{\begin{array}{l} v = 0\ \ \mathrm{m\,s^{-1}} \\ a = -A\omega^2 = -0,25\ \mathrm{m\,s^{-2}} \end{array}\right.$

o fer les derivades de la posició.

% A l'original aquesta nota és al final del problema, després de l'apartat b, però es refereix a l'apartat a.
\textbf{(*)} l'enunciat podria prestar-se a una doble lectura i per tant es donarà com a vàlid interpretar que el període és 6,28~s.
\end{solapartat}
